% !TEX root = ../../draft_v2.tex
%==============================================================================
%  Appendix D7 : The disagreement-node tender game -- deriving lambda
%  ----------------------------------------------------------------------------
%  Drop-in appendix for "Liquidity, Activism Disclosure, and Takeover Premia".
%  Uses the body macros \E, \PP, \1, \mtil, \Dtil and biblatex \citet/\citep.
%
%  WHAT THIS RECORD DOES.  Appendix D4 (= draft Appendix app:bg) honestly
%  relabelled the threat-point shift as Condition BG after five attempts to
%  derive it failed the circularity test: asserting lambda>0 as a Lemma merely
%  relabels the assumption.  D4's own diagnosis of the obstruction (draft
%  l.2440-41): "the sign and magnitude of the shift are governed by a
%  pivotality/appropriability coefficient lambda that can only be pinned by a
%  complete tender-game equilibrium."  THIS RECORD SOLVES THAT TENDER GAME.
%
%  The game is the disagreement-node continuation that D4's free-rider
%  heuristic (app:bg-threat) sketches but does not solve: bargaining with the
%  incumbent bidder has broken down; a fringe raider may mount a Grossman-Hart
%  tender offer against the (possibly engagement-improved) freestanding target.
%  Solving it delivers
%
%        lambda  =  1 - q (1-gamma) psi   in [0,1],
%
%  where q = H(phi) is the equilibrium fringe-raid probability (H = entry-cost
%  cdf, phi = charter dilution), gamma in [0,1] is the PORTABILITY of the
%  activist's improvement under a fringe acquirer (complements vs substitutes),
%  and psi in [0,1] is a pivotality factor (psi < 1 when the bloc can block).
%
%  WHY THIS IS NOT THE CIRCULARITY D4 FLAGGED.  lambda is now an equilibrium
%  OUTPUT.  The structural residue moves one altitude level down: from a
%  condition on an endogenous bargaining object (the threat point d(1)-d(0))
%  to primitives of the disagreement-node environment (H, G, phi, gamma,
%  alpha, tau_c).  A3's failure region is now CHARACTERIZED rather than
%  excluded by assumption: lambda = 0 iff a fringe raid is certain (q = 1),
%  fully supersedes the activist's improvement (gamma = 0), and cannot be
%  blocked or blocking protects nothing.  That boundary is economically
%  meaningful (it is the Albuquerque-Fos-Schroth-flavoured region; see
%  Proposition d7:afs) and empirically indexable (Remark d7:testable).
%
%  EPISTEMIC LEDGER (mirrors D4's discipline):
%   * PROVED: the free-rider floor (Lemma d7:floor); entry rule q = H(phi)
%     independent of the engagement state (Lemma d7:entry); the bloc's
%     blocking/acquiescence rule (Lemma d7:bloc); the threat-point shift and
%     the closed form for lambda (Proposition d7:lambda); Condition BG's
%     functional form as a COROLLARY (Corollary d7:BG); A3 from primitives
%     with sharp iff boundary (Theorem d7:A3); exactness of the rho-fold
%     (Remark d7:fold); the PE measured-premium reversal (Proposition d7:afs).
%   * INSTITUTIONAL PRIMITIVES (transparent, discussed, not derived): uniform
%     equal-treatment offers; the Grossman-Hart tie-breaking selection in the
%     tendering subgame; exogenous charter objects (phi, tau_c); one-shot
%     fringe arrival; engagement state a observable at the node.
%   * OPEN (honestly): bidder competition between incumbent and fringe (an
%     auction in the spirit of Bulow-Huang-Klemperer 1999); dynamic/repeated
%     fringe arrival; endogenous charter design of (phi, tau_c).
%==============================================================================

% --- local guards (must precede first use; no-ops if defined in preamble) ----
\providecommand{\Gbar}{\overline{G}}
\makeatletter
\@ifundefined{c@condition}{\theoremstyle{plain}\newtheorem{condition}{Condition}}{}
\makeatother

\section{The disagreement-node tender game: deriving the appropriability coefficient}
\label{app:d7}

This appendix replaces the \emph{stated} threat-point condition of Appendix~\ref{app:bg} (Condition~\ref{cond:bg}) with a \emph{solved game}. Appendix~\ref{app:bg} was deliberate that its load-bearing object --- the engagement-induced shift $d(1)-d(0)=\lambda\,\Dtil_{\mathrm{eng}}$ in the bloc's disagreement payoff --- could "only be pinned by a complete tender-game equilibrium," and carried $\lambda\in(0,1]$ as a reduced-form pivotality parameter with the free-rider logic of \citet{GrossmanHart1980} offered as motivation, not proof (Remark~\ref{rem:bg-threat-honest}). Here we specify and solve that tender game at the disagreement node. The appropriability coefficient emerges in closed form,
\[
\lambda \;=\; 1-q\,(1-\gamma)\,\psi \;\in\;[0,1],
\]
with $q$ the equilibrium probability of a fringe raid, $\gamma$ the portability of the activist's improvement under a fringe acquirer, and $\psi$ a pivotality factor. Condition~\ref{cond:bg} becomes a corollary (its functional form is \emph{derived}), and Assumption~\textup{(A3)} becomes a theorem whose failure region is characterized by primitives rather than excluded by assumption. The honest residue is stated in \S\ref{app:d7-ledger}: what remains assumed are \emph{institutional} primitives of the tender environment (equal-treatment offers, the standard tie-breaking selection, charter parameters), not any condition on the endogenous threat point itself.

\paragraph{Notation.}
Throughout this appendix $\alpha\in(0,1)$ denotes the control bloc's stake and $\tau_{c}\in(0,1]$ the control threshold (fraction of all shares needed for control); these are local to the tender game.\footnote{The symbol $\alpha$ is unrelated to the completion-probability robustness device $\alpha(P,D)$ mentioned in passing in the body; $\tau_{c}$ is unrelated to the discounting horizon $\tau$. We write $\varphi$ for the standard normal density (as in Appendix~\ref{app:bg}) and reserve $\phi$ for the dilution parameter below, flagging the departure from the body's occasional use of $\phi(\cdot)$ for the density.} We write $\Gbar(x)\equiv\PP(S_{F}\ge x)$ for the survival function of the fringe synergy distribution $G$, assumed atomless on its support and with finite mean. We maintain throughout that
\[
\alpha\;<\;\tau_{c}:
\]
the bloc \emph{alone} cannot deliver control (satisfied at the baseline calibration, $\alpha=0.30<\tau_{c}=0.75$). Without it the raider's cheapest successful coalition is a bloc-only lowball at the bloc's reservation, the bloc is held at $y(a)$ in \emph{every} raid, and $\lambda=1$ identically --- the closed form below would not apply (see \S\ref{app:d7-ledger}).

\subsection{Environment}
\label{app:d7-env}

Condition on the disagreement node of Appendix~\ref{app:bg}: negotiation with the incumbent bidder has broken down, and the target is a freestanding firm. As in Appendix~\ref{app:bg}, $a\in\{0,1\}$ indexes the engagement state at the node, and we work conditional on the \emph{realized} engagement improvement $\delta_{e}\ge 0$ (so $\delta_{e}=\Delta_{\mathrm{eng}}$ if the campaign succeeded, $\delta_{e}=0$ otherwise); the fold back to expected objects $\Dtil_{\mathrm{eng}}=\rho\Delta_{\mathrm{eng}}$ is exact and is performed in Remark~\ref{rem:d7-fold}, consistent with the two-draw event tree of Lemma~\ref{lem:bg-tree}. The firm's standalone per-share value at the node is
\[
y(a)\;=\;w+a\,\delta_{e}.
\]

\textbf{Players.} (i) A \emph{fringe raider} with privately observed type $(c_{F},S_{F})\sim H\otimes G$, independent draws: $c_{F}\ge0$ is the cost of mounting a tender offer, and $S_{F}\ge 0$ is the raider's synergy over the firm \emph{as it stands under the raider's plan} (below). (ii) The \emph{control bloc}, holding stake $\alpha\in(0,1)$, strategic. (iii) A dispersed \emph{float} of measure $1-\alpha$, atomistic.

\textbf{Portability.} If the raider takes control, the post-takeover per-share gross value is
\[
Z(a)\;=\;w+\gamma\,a\,\delta_{e}+S_{F},
\qquad \gamma\in[0,1].
\]
The parameter $\gamma$ is the fraction of the activist's improvement that \emph{survives a change of control}: $\gamma=1$ when the improvement is portable (governance fixes, balance-sheet repairs, divestitures already executed --- complements to any owner), $\gamma=0$ when the raider's plan supersedes it (the raider re-optimizes operations wholesale, so the activist's program is redundant --- substitutes). Intermediate $\gamma$ mixes the two.

\textbf{Dilution.} As in \citet{GrossmanHart1980}, the corporate charter permits a controlling raider to dilute minority (non-tendered) shares by $\phi\ge0$ per share; the diverted value accrues to the raider.

\textbf{Timing.}
\begin{enumerate}
\item The raider observes $(c_{F},S_{F})$ and the public state $a$, and chooses whether to \emph{enter} (sinking $c_{F}$); upon entry it posts a \emph{uniform, conditional, unrestricted} tender offer at price $\hat b$ per share, conditional on receiving at least $\tau_{c}$ of the shares.
\item Float members tender atomistically; the bloc tenders strategically (all-or-none is without loss under a uniform price).
\item If the tendered mass is at least $\tau_{c}$, control transfers: tendered shares receive $\hat b$, non-tendered shares are worth $Z(a)-\phi$. Otherwise the offer lapses and the status quo $y(a)$ prevails.
\end{enumerate}

\textbf{Equilibrium concept.} Subgame-perfect equilibrium of the entry-offer-tender game, with the standard tie-breaking selection in the tendering subgame: shareholders who are exactly indifferent tender. This is the selection under which value-improving raids are possible at all in the atomistic limit \citep{GrossmanHart1980}; \citet{BagnoliLipman1988} provide the finite-shareholder foundation in which pivotal tendering delivers this outcome as the limit of exact equilibria, and \citet{HolmstromNalebuff1992} characterize the mixed equilibria we do not select. We say the bloc is \emph{pivotal} if the float alone cannot deliver control, i.e.\ $1-\alpha<\tau_{c}$.

\subsection{The tendering subgame and the free-rider floor}

\begin{lemma}[Optimal offers: the free-rider floor and the reservation top-up]
\label{lem:d7-floor}
Fix an entered raider and the state $a$, and define the \emph{acquiescence region} as all of the support of $S_F$ when the bloc is non-pivotal, and $\{S_{F}\ge\phi+(1-\gamma)a\delta_{e}\}$ when it is pivotal (Lemma~\ref{lem:d7-bloc}); its complement (pivotal only) is the \emph{blocking region}. In every subgame-perfect tendering outcome under the stated selection:
\begin{enumerate}
\item[\textup{(i)}] On the acquiescence region the optimal offer is the \emph{free-rider floor}
\[
\hat b^{\ast}(a)\;=\;Z(a)-\phi ,
\]
at which the entire float tenders. The raider's payoff per share it does \emph{not} acquire is the dilution $\phi$, and per share it \emph{does} acquire is $Z(a)-\hat b^{\ast}(a)=\phi$; its gross payoff from a successful floor raid is $\phi$ on the full register, independently of the tendered mass, and net of cost $\phi-c_{F}$.
\item[\textup{(ii)}] On the blocking region the floor offer fails (the pivotal bloc refuses), and the raider's best successful offer is the \emph{reservation top-up} $\hat b=y(a)>Z(a)-\phi$: the float strictly tenders, the indifferent bloc tenders under the selection, the raider acquires the full register and nets
\[
Z(a)-y(a)-c_{F}\;=\;S_{F}-(1-\gamma)\,a\,\delta_{e}-c_{F},
\]
undertaking the raid iff this is nonnegative. In every top-up raid the bloc is paid \emph{exactly} its reservation $y(a)$.
\end{enumerate}
\end{lemma}

\begin{proof}
An atomistic float member is never pivotal, so conditional on the raid succeeding her choice is $\hat b$ (tender) against $Z(a)-\phi$ (retain a diluted share); conditional on failure the conditional offer lapses and both choices give $y(a)$. She therefore tenders iff $\hat b\ge Z(a)-\phi$, with the selection resolving indifference toward tendering.

(i) On the acquiescence region the floor succeeds: the float tenders, and (pivotal case) the bloc tenders by Lemma~\ref{lem:d7-bloc}(ii). Any $\hat b<Z(a)-\phi$ loses the entire float, and with $\alpha<\tau_{c}$ bloc shares alone cannot deliver control, so the raid fails. Any $\hat b>Z(a)-\phi$ keeps the success event unchanged while strictly raising the price on every tendered share, hence is dominated. At the floor: each acquired share costs $Z(a)-\phi$ and is worth $Z(a)$, margin $\phi$; each non-acquired share is diluted by $\phi$, accruing to the raider; summing over the unit register gives $\phi$ gross, $\phi-c_{F}$ net, independent of the tendered mass and of $a$.

(ii) On the blocking region $y(a)>Z(a)-\phi$. Offers $\hat b\in[Z(a)-\phi,\,y(a))$ attract the float but not the bloc, delivering mass $1-\alpha<\tau_{c}$: failure. The top-up $\hat b=y(a)$ exceeds the float's reservation (strict tendering) and meets the bloc's (indifference resolved toward tendering): mass one, success. Raising $\hat b$ above $y(a)$ only pays more on every share. At $\hat b=y(a)$ the raider buys the whole register at $y(a)$ per share, each worth $Z(a)$ under its control, with no minority left to dilute: net payoff $Z(a)-y(a)-c_{F}=S_{F}-(1-\gamma)a\delta_{e}-c_{F}$. Entry on the blocking region therefore occurs exactly when this is nonnegative; the bloc, tendering at its reservation, receives $y(a)$.
\end{proof}

The floor is the free-rider result: the raider's synergy $S_{F}$ passes \emph{entirely} to the shareholders through the offer price; the raider's only rent on floor raids is the charter dilution $\phi$. Top-up raids exist precisely because a reluctant pivotal bloc can be bought out at its outside option --- but for that very reason they transfer \emph{no} surplus to the bloc, which is what keeps the threat point below clean of them.

\subsection{Entry}

\begin{lemma}[Floor raids are dilution-financed and state-blind]
\label{lem:d7-entry}
Floor raids --- raids at the offer $\hat b^{\ast}(a)=Z(a)-\phi$, which are the only raids in which the bloc earns more than its reservation --- occur iff $c_{F}\le\phi$ (and, if the bloc is pivotal, the acquiescence event of Lemma~\ref{lem:d7-bloc} holds). The floor-feasibility event $\{c_{F}\le\phi\}$ has probability
\[
q\;\equiv\;H(\phi)\;\in\;[0,1],
\]
independent of the engagement state $a$ and of the synergy draw $S_{F}$. Top-up raids (Lemma~\ref{lem:d7-floor}(ii)) occur on a state-dependent set --- engagement \emph{shrinks} it --- but pay the bloc exactly $y(a)$ and therefore do not enter the threat-point shift.
\end{lemma}

\begin{proof}
By Lemma~\ref{lem:d7-floor}(i) a floor raid pays $\phi-c_{F}$ regardless of $(a,S_{F})$ and of the tendered mass; not entering pays $0$ (the offer is conditional). On the acquiescence region the floor is the raider's best offer, so floor raids occur exactly on $\{c_{F}\le\phi\}$ intersected with acquiescence, and $\PP(c_{F}\le\phi)=H(\phi)$ with independence from $(a,S_{F})$ immediate. On the blocking region the optimal offer is the top-up of Lemma~\ref{lem:d7-floor}(ii), profitable on $\{S_{F}-(1-\gamma)a\delta_{e}\ge c_{F}\}$ --- a set decreasing in $a\delta_{e}$ --- in which the bloc receives its reservation $y(a)$ exactly.
\end{proof}

The economically important part is that the \emph{bloc-surplus-relevant} raid margin is state-blind: on floor raids the raider's profit is $\phi-c_{F}$ whatever the engagement state, so engagement affects the bloc's takeover option only through its own willingness to part with an improved firm --- not through raider selection. Total raid incidence \emph{is} state-dependent (engagement deters some top-up raids), but those raids hold the bloc at its outside option and so contribute nothing to the shift $d(1)-d(0)$. This is what keeps the threat-point shift below clean of entry-selection terms.

\subsection{The bloc at the node: blocking and acquiescence}

\begin{lemma}[Blocking rule]
\label{lem:d7-bloc}
Consider an entered raider offering the floor $\hat b^{\ast}(a)=Z(a)-\phi$.
\begin{enumerate}
\item[\textup{(i)}] \textbf{Non-pivotal bloc} ($1-\alpha\ge\tau_{c}$): the raid succeeds for every $S_{F}$; the bloc's per-share payoff is $Z(a)-\phi$ whether it tenders or retains.
\item[\textup{(ii)}] \textbf{Pivotal bloc} ($1-\alpha<\tau_{c}$): the raid succeeds iff the bloc tenders, and the bloc tenders iff
\[
\hat b^{\ast}(a)\;\ge\;y(a)
\quad\Longleftrightarrow\quad
S_{F}\;\ge\;\phi+(1-\gamma)\,a\,\delta_{e}.
\]
On the acquiescence set the bloc's per-share payoff is $Z(a)-\phi$; off it, the \emph{floor} raid fails, and the bloc receives $y(a)$ --- either because no raid occurs, or because the raider switches to the reservation top-up of Lemma~\ref{lem:d7-floor}(ii), which pays the bloc exactly $y(a)$.
\end{enumerate}
\end{lemma}

\begin{proof}
(i) With $1-\alpha\ge\tau_{c}$ the float alone delivers control (Lemma~\ref{lem:d7-floor}(i): all of it tenders at the floor). The bloc tendering yields $\hat b^{\ast}=Z(a)-\phi$; retaining yields the diluted value $Z(a)-\phi$. Identical.
(ii) With $1-\alpha<\tau_{c}$ the success event requires bloc shares. If the bloc tenders at the floor, it receives $\hat b^{\ast}(a)$; if it refuses, the conditional floor offer lapses and it retains a share worth $y(a)$ (or is bought out at $y(a)$ in a top-up raid). It tenders at the floor iff $Z(a)-\phi\ge y(a)$, i.e.\ $w+\gamma a\delta_{e}+S_{F}-\phi\ge w+a\delta_{e}$, i.e.\ $S_{F}\ge\phi+(1-\gamma)a\delta_{e}$. (Indifference, an event of $G$-measure zero, is resolved toward tendering.) Off the acquiescence set the raider's optimal response is characterized by Lemma~\ref{lem:d7-floor}(ii): the top-up succeeds when profitable, and in it the bloc --- tendering at its reservation under the selection --- receives exactly $y(a)$. In every branch off the acquiescence set the bloc's payoff is therefore $y(a)$.
\end{proof}

\subsection{The threat point and the appropriability coefficient}

\begin{proposition}[The threat-point shift, and $\lambda$ in closed form]
\label{prop:d7-lambda}
Let $d(a)$ be the bloc's per-share equilibrium continuation value at the disagreement node, conditional on the realized improvement $\delta_{e}$. Then
\begin{align}
\text{non-pivotal:}\quad d(a)&\;=\;w+a\delta_{e}+q\,\bigl(\E[S_{F}]-\phi-(1-\gamma)a\delta_{e}\bigr),
\label{eq:d7-dNP}\\
\text{pivotal:}\quad d(a)&\;=\;w+a\delta_{e}+q\,\E\bigl[(S_{F}-\phi-(1-\gamma)a\delta_{e})^{+}\bigr].
\label{eq:d7-dP}
\end{align}
In both regimes the engagement-induced shift is linear in the realized improvement,
\[
d(1)-d(0)\;=\;\lambda\,\delta_{e},
\qquad
\boxed{\;\lambda\;=\;1-q\,(1-\gamma)\,\psi\;\in\;[0,1],\;}
\]
where $\psi=1$ in the non-pivotal regime and, in the pivotal regime with $\gamma<1$ and $\delta_{e}>0$,
\[
\psi\;=\;\frac{1}{(1-\gamma)\,\delta_{e}}\int_{0}^{(1-\gamma)\delta_{e}}\Gbar(\phi+t)\,dt\;\in\;[0,1]
\]
(the average probability, along the forgone-improvement interval, that the fringe synergy clears the bloc's reservation); when $\gamma=1$ or $\delta_{e}=0$ the $\psi$-term is multiplied by zero and $\lambda=1$. Moreover $\psi\le 1$ implies $\lambda^{\textup{pivotal}}\ge\lambda^{\textup{non-pivotal}}$ at common $(q,\gamma)$: pivotality protects the threat point.
\end{proposition}

\begin{proof}
\emph{Non-pivotal.} By Lemmas~\ref{lem:d7-entry}--\ref{lem:d7-bloc}(i), with probability $q$ a raider enters and the raid succeeds for every $S_{F}$, paying the bloc $Z(a)-\phi=w+\gamma a\delta_{e}+S_{F}-\phi$; with probability $1-q$ the status quo $y(a)=w+a\delta_{e}$ prevails. Taking expectations over $S_{F}$ (independent of entry) gives
\[
d(a)=(1-q)(w+a\delta_{e})+q\,(w+\gamma a\delta_{e}+\E[S_{F}]-\phi),
\]
which rearranges to \eqref{eq:d7-dNP}. Differencing: $d(1)-d(0)=(1-q)\delta_{e}+q\gamma\delta_{e}=\delta_{e}\,[1-q(1-\gamma)]$.

\emph{Pivotal.} By Lemmas~\ref{lem:d7-floor}--\ref{lem:d7-bloc}, the bloc earns more than its reservation only in \emph{floor} raids, which occur on $\{c_{F}\le\phi\}\cap\{S_{F}\ge\phi+(1-\gamma)a\delta_{e}\}$ and pay it $Z(a)-\phi$; in every other branch --- no entry, failed floor, or a top-up raid that buys the bloc out at its outside option --- the bloc has exactly $y(a)$ (the outside-option principle: marginal raids extract the bloc's full surplus). Hence
\[
d(a)=y(a)+q\,\E\bigl[\bigl(Z(a)-\phi-y(a)\bigr)\,\1\{S_{F}\ge\phi+(1-\gamma)a\delta_{e}\}\bigr]
=y(a)+q\,\E\bigl[(S_{F}-\phi-(1-\gamma)a\delta_{e})^{+}\bigr],
\]
using $Z(a)-\phi-y(a)=S_{F}-\phi-(1-\gamma)a\delta_{e}$, which is exactly the positive part on the acquiescence set. This is \eqref{eq:d7-dP}. Differencing, with $X\equiv S_{F}-\phi$ and $c\equiv(1-\gamma)\delta_{e}\ge0$:
\[
d(1)-d(0)=\delta_{e}+q\,\bigl(\E[(X-c)^{+}]-\E[X^{+}]\bigr)
=\delta_{e}-q\int_{0}^{c}\PP(X\ge t)\,dt ,
\]
where the last step is the layer-cake identity $\E[X^{+}]-\E[(X-c)^{+}]=\int_{0}^{c}\PP(X\ge t)\,dt$ (Fubini on $\E\int_{0}^{\infty}\1\{t\le X\}dt$, valid for any $X$ with $\E|X|<\infty$ and $c\ge0$). Substituting $\PP(X\ge t)=\Gbar(\phi+t)$ and factoring out $c=(1-\gamma)\delta_{e}$ yields $d(1)-d(0)=\delta_{e}[1-q(1-\gamma)\psi]$ with $\psi$ as displayed. Bounds: $0\le\Gbar\le1$ gives $\psi\in[0,1]$, and $q,(1-\gamma),\psi\in[0,1]$ give $\lambda\in[0,1]$. The regime comparison is immediate from $\psi\le1$.
\end{proof}

\begin{corollary}[Condition BG's form is derived, with a microfounded level]
\label{cor:d7-BG}
Define $w'\equiv d(0)$ from \eqref{eq:d7-dNP}--\eqref{eq:d7-dP} (the standalone continuation value absent engagement, now inclusive of the fringe option value) and $\lambda$ as in Proposition~\ref{prop:d7-lambda}. Then the bloc's disagreement payoff satisfies \emph{exactly}
\[
d(a)\;=\;w'+a\,\lambda\,\delta_{e},
\]
i.e.\ equation~\eqref{eq:bg-disagree} of Condition~\ref{cond:bg} holds with $\lambda=1-q(1-\gamma)\psi$ an equilibrium object. Every downstream result of Appendix~\ref{app:bg} that takes Condition~\ref{cond:bg} as input (Lemma~\ref{lem:bg-split}, the wedge \eqref{eq:bg-wedge}, Proposition~\ref{prop:bg-reduced}, Theorem~\ref{thm:bg-A3}) now operates on a derived premise, with one caveat: $\lambda>0$ is no longer guaranteed by assumption --- it holds exactly on the region characterized in Theorem~\ref{thm:d7-A3}, and the wedge formula extends continuously to $\lambda=0$.
\end{corollary}

\begin{proof}
Linearity of $d(a)$ in $a$ with slope $\lambda\delta_{e}$ is the content of Proposition~\ref{prop:d7-lambda}; matching coefficients against \eqref{eq:bg-disagree} (with $\delta_{e}$ in place of $\Dtil_{\mathrm{eng}}$ pre-fold; see Remark~\ref{rem:d7-fold}) gives the display. The downstream results of Appendix~\ref{app:bg} use only the functional form and $\lambda\ge0$; inspection of their proofs shows none uses $\lambda>0$ except where flagged.
\end{proof}

\begin{remark}[The fold back to $\Dtil_{\mathrm{eng}}$ is exact]
\label{rem:d7-fold}
Proposition~\ref{prop:d7-lambda} conditions on the realized improvement $\delta_{e}\in\{0,\Delta_{\mathrm{eng}}\}$. Taking expectations over the campaign-success draw $B_{1}\sim\mathrm{Bern}(\rho)$ of Lemma~\ref{lem:bg-tree}: the shift is $0$ when $B_{1}=0$ and $\lambda(\Delta_{\mathrm{eng}})\,\Delta_{\mathrm{eng}}$ when $B_{1}=1$, so the expected shift is $\rho\,\lambda\,\Delta_{\mathrm{eng}}=\lambda\,\Dtil_{\mathrm{eng}}$ with $\lambda\equiv\lambda(\Delta_{\mathrm{eng}})$ evaluated at the realized improvement. No approximation is involved, and the second $\rho$ of the body's fold (premium-state arrival $B_{2}$) composes exactly as in Lemma~\ref{lem:bg-tree}, preserving $\mtil-m_{0}=\rho^{2}(1-\theta)\lambda\Delta_{\mathrm{eng}}$. Note one refinement relative to the reduced form: $\lambda(\cdot)$ depends on the \emph{size} of the realized improvement in the pivotal regime (through $\psi$), a state-dependence the constant-$\lambda$ reduced form suppresses; quantitative statements therefore hold $\lambda$ at $\lambda(\Delta_{\mathrm{eng}})$.
\end{remark}

\subsection{Assumption (A3) from primitives}

\begin{theorem}[A3 holds except on a characterized boundary]
\label{thm:d7-A3}
Maintain $\theta<1$ and $\rho\,\Delta_{\mathrm{eng}}>0$, and let $\lambda=1-q(1-\gamma)\psi$ from Proposition~\ref{prop:d7-lambda}. Then the bargained premia of Appendix~\ref{app:bg} satisfy
\[
\mtil-m_{0}\;=\;\rho^{2}(1-\theta)\,\lambda\,\Delta_{\mathrm{eng}}\;>\;0
\qquad\Longleftrightarrow\qquad
\lambda>0,
\]
and $\lambda=0$ --- equivalently, Assumption~\textup{(A3)} fails --- if and only if \emph{all three} of the following hold:
\begin{enumerate}
\item[\textup{(F1)}] fringe raids are \emph{certain}: $q=H(\phi)=1$;
\item[\textup{(F2)}] the activist's improvement is \emph{fully superseded} by the raider's plan: $\gamma=0$;
\item[\textup{(F3)}] blocking is \emph{unavailable or worthless}: the bloc is non-pivotal ($1-\alpha\ge\tau_{c}$), \emph{or} it is pivotal but the fringe synergy clears the bloc's reservation with probability one, $\PP\bigl(S_{F}\ge\phi+\Delta_{\mathrm{eng}}\bigr)=1$.
\end{enumerate}
In particular each of the following alone implies $\lambda>0$: some raids fail to materialize ($q<1$); any portability ($\gamma>0$); a pivotal bloc facing a fringe whose synergy falls short of $\phi+\Delta_{\mathrm{eng}}$ with positive probability.
\end{theorem}

\begin{proof}
The equivalence $\mtil-m_{0}>0\Leftrightarrow\lambda>0$ is Theorem~\ref{thm:bg-A3}'s display with $\lambda$ substituted from Corollary~\ref{cor:d7-BG}, the other factors being positive by hypothesis. It remains to characterize $\lambda=0$, i.e.\ $q(1-\gamma)\psi=1$. Since each factor lies in $[0,1]$, the product equals one iff every factor equals one. $q=1$ is (F1); $1-\gamma=1$ is (F2). Given $\gamma=0$: in the non-pivotal regime $\psi\equiv1$, which is the first disjunct of (F3); in the pivotal regime $\psi=1$ iff $\int_{0}^{\delta_{e}}\Gbar(\phi+t)\,dt=\delta_{e}$ iff $\Gbar(\phi+t)=1$ for a.e.\ $t\in(0,\delta_{e})$; since $\Gbar$ is nonincreasing and right-continuous this holds iff $\Gbar(\phi+\delta_{e}^{-})=1$, i.e.\ $\PP(S_{F}\ge\phi+\delta_{e})=1$ for $\delta_{e}=\Delta_{\mathrm{eng}}$ by atomlessness of $G$ --- the second disjunct of (F3). The "in particular" statement negates each factor.
\end{proof}

\begin{remark}[Comparative statics of appropriability]
\label{rem:d7-compstat}
From $\lambda=1-q(1-\gamma)\psi$ with $q=H(\phi)$:
(i) $\lambda$ is nonincreasing in fringe-market strength: $\partial\lambda/\partial q=-(1-\gamma)\psi\le0$; a hotter takeover market \emph{erodes} the activist's bargaining-relevant appropriability even as it raises the level of the bloc's outside option.
(ii) $\lambda$ is nondecreasing in portability: in the pivotal regime $(1-\gamma)\psi=\frac{1}{\delta_{e}}\int_{0}^{(1-\gamma)\delta_{e}}\Gbar(\phi+t)\,dt$, whose $\gamma$-derivative is $-\Gbar(\phi+(1-\gamma)\delta_{e})\le0$, so $\partial\lambda/\partial\gamma=q\,\Gbar(\phi+(1-\gamma)\delta_{e})\ge0$; in the non-pivotal regime $\partial\lambda/\partial\gamma=q\ge0$.
(iii) Pivotality is protective and discontinuous: as $\alpha$ crosses $1-\tau_{c}$ from below, $\psi$ drops from $1$ to the integral average, so $\lambda$ jumps \emph{up} by $q(1-\gamma)\bigl(1-\psi\bigr)\ge0$ --- a testable stake-size threshold effect.
(iv) Dilution $\phi$ is two-edged in the pivotal regime: $\partial\lambda/\partial\phi=-(1-\gamma)\psi\,H'(\phi)+q\,\frac{1}{\delta_{e}}\bigl[\Gbar(\phi)-\Gbar(\phi+(1-\gamma)\delta_{e})\bigr]\cdot$ $(\text{sign depends on } H' \text{ vs the } G\text{-mass on }(\phi,\phi+(1-\gamma)\delta_{e}])$: easier dilution invites more raids ($q\uparrow$, eroding $\lambda$) but also lowers the floor relative to the bloc's reservation, thinning the acquiescence set ($\psi\downarrow$, protecting $\lambda$). We leave $\phi$ unsigned and verify both directions numerically.
\end{remark}

\subsection{Measured premia: rationalizing the estimated negative activism--premium effect}
\label{app:d7-afs}

\citet{CelentanoLevine2025} estimate that bid premia are 13.7\% (5.2pp) lower when an activist is present in takeover negotiations, relative to a no-activist counterfactual.\footnote{In their structural model the channel is bargaining: the activist lowers the board's entrenchment cost, which lowers the price the board demands. The mechanism developed here is complementary and differently indexed --- see Remark~\ref{rem:d7-afs-comp}. \citet{AlbuquerqueFosSchroth2022} contain no takeover-premium estimate; their role in this section is different and essential: they quantify the \emph{capitalization} of activism in prices at disclosure (a 6.34\% announcement return, roughly three-quarters of it anticipated treatment), which is precisely the denominator effect the proposition below exploits.} At face value a negative activism--premium relation contradicts $m_{1}>m_{0}$. The model says otherwise: the empirical object is not the wedge. Bid premia are measured against a market price that \emph{already capitalizes} activism beliefs --- the body's Proposition~\ref{prop:price-decomp} prices in both the expected standalone improvement $\Dtil\,\pi(X,D)$ and the activism premium $(\mtil-m_{0})\pi(X,D)$ --- so activism moves the \emph{denominator} as well as the numerator. When appropriability is low, the denominator wins.

\begin{proposition}[Measured-premium reversal under low appropriability]
\label{prop:d7-afs}
Fix the blockholder's cutoffs and an information state $(X,D)$ as in Proposition~\ref{prop:disclosure-attenuation}, and let $P(\pi)$ denote the pricing fixed point of equation~\eqref{eq:pricing} viewed as a function of the engagement posterior $\pi=\pi(X,D)$, with $\bar m(\pi)=m_{0}+(\mtil-m_{0})\pi$ and bid probability $p(P,\pi)$ from \eqref{eq:bid-prob}. Define the \emph{measured premium}
\[
M(\pi)\;\equiv\;\frac{\bar m(\pi)}{P(\pi)} ,
\]
the expected offer over the prevailing market price, net of one. Maintain Assumption~\textup{(A5a)} and $m_{0}>0$. Then
\[
\operatorname{sgn} M'(\pi)\;=\;\operatorname{sgn}\Bigl[(\mtil-m_{0})\,P(\pi)\;-\;\bar m(\pi)\,P'(\pi)\Bigr],
\]
and in the $\lambda\to0$ limit $P'(\pi)>0$ whenever $\Dtil>0$.\footnote{At general $\lambda$ the sign of $P'$ is ambiguous: the deterrence term $p_{\pi}\,(P+\bar m-\E[v|X,D]-\Dtil\pi)$ can dominate $(1-p)\Dtil$ when the wedge is large relative to $\Dtil$, so the clause is claimed only in the limit, which is all the conclusion below requires. A sufficient condition at general $\lambda$ is $p\,(\mtil-m_{0})+(1-p)\,\Dtil>\varphi(t)\,(\mtil-m_{0})\,(P+\bar m-\E[v|X,D]-\Dtil\pi)/\sigma_{\xi}$.}
Since $\mtil-m_{0}=\rho^{2}(1-\theta)\lambda\Delta_{\mathrm{eng}}\to0$ as $\lambda\to0$ while $\bar m(\pi)\to m_{0}>0$ and
\[
P'(\pi)\;\longrightarrow\;\frac{\delta\,(1-p)\,\Dtil}{1-\delta p-\delta\,\frac{\partial p}{\partial P}\,\bigl(P+m_{0}-\E[v\mid X,D]-\Dtil\pi\bigr)}\;>\;0
\qquad(\lambda\to0),
\]
there exists $\lambda_{\mathrm{crit}}>0$ such that for all $\lambda<\lambda_{\mathrm{crit}}$ the measured premium is \emph{strictly decreasing} in activism evidence, $M'(\pi)<0$ for all interior $\pi$ --- even though $m_{1}\ge m_{0}$ throughout. Activism then lowers measured bid premia as \citet{CelentanoLevine2025} estimate, through the price channel, with the wedge's sign intact.
\end{proposition}

\begin{proof}
Differentiate $M=\bar m/P$: $M'=[(\mtil-m_{0})P-\bar m P']/P^{2}$, giving the sign display. For $P'$: the fixed point is $P=\delta[\,p(P,\pi)\,(P+\bar m(\pi))+(1-p(P,\pi))\,(\E[v\mid X,D]+\Dtil\pi)\,]$. The implicit function theorem applies: the denominator $1-\delta p-\delta\,p_{P}\,(P+\bar m-\E[v|X,D]-\Dtil\pi)$ is bounded below by $1-\delta\bigl[p+|P+\bar m-\E[v|X,D]-\Dtil\pi|\,\varphi(0)/\sigma_{\xi}\bigr]>0$ under the contraction bound of Assumption~\textup{(A5a)} (see also Remark~\ref{rem:A5margins} in the body: at the baseline calibration the conservative sufficient form of A5a is not met, so the numerical $\lambda_{\mathrm{crit}}$ below rests on the solver's observed contraction rather than the maintained bound), and it gives
\[
P'(\pi)=\frac{\delta\bigl[p\,(\mtil-m_{0})+(1-p)\,\Dtil+p_{\pi}\,\bigl(P+\bar m-\E[v|X,D]-\Dtil\pi\bigr)\bigr]}{1-\delta p-\delta\,p_{P}\,\bigl(P+\bar m-\E[v|X,D]-\Dtil\pi\bigr)} .
\]
As $\lambda\to0$: $\mtil-m_{0}\to0$ directly, and $p_{\pi}=-\varphi(\cdot)\,(\mtil-m_{0})/\sigma_{\xi}\to0$ since $\pi$ enters $p$ only through $\bar m$; the displayed limit follows, and it is strictly positive for $\Dtil>0$, $p<1$. All equilibrium objects at fixed cutoffs are continuous in $\lambda$ (they are compositions of $\Phi$, $\varphi$, and the fixed point, which is continuous by A5 and the uniform contraction). Hence $M'\to-m_{0}P'/P^{2}<0$ uniformly on compact interior $\pi$-sets, and continuity delivers $\lambda_{\mathrm{crit}}>0$.
\end{proof}

\begin{remark}[Composition reinforces the reversal; what would distinguish the stories]
\label{rem:d7-afs-comp}
Two further forces push the same way in the cross-section. First, deterrence: $\partial p/\partial\pi<0$ when $\mtil>m_{0}$ (body, \S Bid Incidence), so high-$\pi$ states select toward deals that clear a higher hurdle --- fewer, not richer, marginal deals. Second, heterogeneity in $\lambda$: by Theorem~\ref{thm:d7-A3}, campaigns at targets whose natural acquirers are \emph{superseding} (strategic buyers re-optimizing wholesale: $\gamma\approx0$) in hot deal markets ($q\approx1$) carry $\lambda\approx0$ and hence near-zero wedges; if such targets are over-represented among consummated takeovers, the estimated average effect tilts negative. The model's distinguishing prediction --- against \citet{CelentanoLevine2025}'s board-bargaining channel, in which the premium effect moves only with activist presence: the activism effect on premia should be \emph{less negative} (i) where improvements are portable (financial acquirers, asset redeployability), (ii) where the activist's stake is pivotal relative to control thresholds, and (iii) in cold fringe markets --- because these are exactly the high-$\lambda$ cells. Cross-sections in $(q,\gamma,\psi)$ separate the two mechanisms for the same average sign.
\end{remark}

\begin{remark}[Empirical content of the primitives]
\label{rem:d7-testable}
Unlike the reduced-form $\lambda$, the derived coefficient is indexable: $q$ by fringe M\&A intensity in the target's industry-year; $\gamma$ by the strategic-vs-financial composition of likely acquirers and asset redeployability; $\phi$ by charter/state anti-takeover provisions governing freeze-outs and squeeze-outs; pivotality by the bloc's stake against control thresholds. This converts Remark~\ref{rem:bg-ident}'s non-identification (only the product $(1-\theta)\lambda\rho\Delta_{\mathrm{eng}}$ pinned) into an over-identifying restriction set: variation in $(q,\gamma,\phi,\alpha)$ should move estimated wedges with the signs of Remark~\ref{rem:d7-compstat}.
\end{remark}

\subsection{What is solved, what is institutional, what remains open}
\label{app:d7-ledger}

\emph{Solved (no residual condition on endogenous objects):} the tendering subgame and free-rider floor (Lemma~\ref{lem:d7-floor}); entry (Lemma~\ref{lem:d7-entry}); the blocking rule (Lemma~\ref{lem:d7-bloc}); the threat point, $\lambda=1-q(1-\gamma)\psi$, and the exact recovery of Condition~\ref{cond:bg}'s form (Proposition~\ref{prop:d7-lambda}, Corollary~\ref{cor:d7-BG}); A3's iff boundary (Theorem~\ref{thm:d7-A3}); the fold (Remark~\ref{rem:d7-fold}); the measured-premium reversal at fixed cutoffs (Proposition~\ref{prop:d7-afs}).

\emph{Institutional primitives (transparent and discussed, in the tradition of taking the tender institution as given):} uniform equal-treatment conditional offers (discriminatory offers would let the raider pay the pivotal bloc exactly its reservation in \emph{every} raid while holding the float at the floor; the bloc would then be held at $y(a)$ throughout, giving $\lambda=1$ --- its maximum. Equal treatment thus yields the \emph{smallest} $\lambda$ and is conservative for A3's failure region; it is also the empirically relevant institution under Williams-Act-type equal treatment); the maintained domain restriction $\alpha<\tau_{c}$ (if the bloc alone delivers control, the raider's bloc-only lowball holds it at reservation in every raid and $\lambda=1$ identically --- A3 then holds trivially, so the restriction is conservative as well); the \citet{GrossmanHart1980} indifference-tendering selection, with \citet{BagnoliLipman1988} as its finite-shareholder foundation and \citet{HolmstromNalebuff1992}'s mixed equilibria not selected; exogenous charter objects $(\phi,\tau_{c})$; one-shot fringe arrival with type $(c_{F},S_{F})\sim H\otimes G$ independent; the engagement state observable at the node.

\emph{Open (stated plainly):} competition between the incumbent bidder and the fringe (a takeover auction in the spirit of \citet{BulowHuangKlemperer1999}, in which the engagement stake may also confer a toehold advantage); repeated fringe arrival and option-value timing at the disagreement node; endogenous design of $(\phi,\tau_{c})$; and the economic identification of the second fold draw $B_{2}$ of Lemma~\ref{lem:bg-tree} (``premium-state arrival'') as an event distinct from campaign success --- under a single-event reading the wedge folds once, $\mtil-m_{0}=\rho(1-\theta)\lambda\Delta_{\mathrm{eng}}$, and the calibrated $\Delta_{\mathrm{eng}}$ shifts from $\approx0.516$ to $\approx0.46$ with all signs and the architecture unchanged. None of these touches the \emph{sign} architecture: each would alter the level of $d(a)$ and potentially the magnitude of $\lambda$, but the linear-in-$a$ form of Corollary~\ref{cor:d7-BG} and the boundary logic of Theorem~\ref{thm:d7-A3} (raids must be certain, superseding, and unblockable to kill the wedge) are robust to them in the obvious directions.

% (local guards are at the top of this file, before first use)
